% ============================================================
% 实证论文 LaTeX 模板 — AER/QJE/管理世界 风格
% 使用: 替换 [方括号] 中的内容
% 编译: xelatex → bibtex → xelatex × 2
% ============================================================
\documentclass[12pt,a4paper]{article}

% === 基础包 ===
\usepackage[utf8]{inputenc}
\usepackage[T1]{fontenc}
\usepackage{amsmath,amssymb,amsthm}
\usepackage{graphicx}
\usepackage{booktabs}        % 三线表
\usepackage{tabularx}
\usepackage{multirow}
\usepackage{float}
\usepackage{caption}
\usepackage{subcaption}
\usepackage[colorlinks=true,linkcolor=blue,citecolor=blue,urlcolor=blue]{hyperref}
\usepackage{natbib}          % 引用格式
\usepackage{setspace}
\usepackage{geometry}
\usepackage{threeparttable}  % 表格注释

% === 中文支持（如需要） ===
% \usepackage{xeCJK}
% \setCJKmainfont{SimSun}     % 宋体
% \setCJKmonofont{SimHei}     % 黑体

% === 页面设置 ===
\geometry{left=2.5cm, right=2.5cm, top=2.5cm, bottom=2.5cm}
\onehalfspacing               % 1.5倍行距

% === 标题信息 ===
\title{\textbf{[论文标题]}\thanks{[致谢/基金信息]}}
\author{
    [作者1]\thanks{[单位、邮箱]} \and
    [作者2]\thanks{[单位、邮箱]}
}
\date{\today}

\begin{document}

\maketitle

\begin{abstract}
\noindent
[摘要内容，250-300词。研究问题→方法→核心发现→政策启示]

\medskip
\noindent\textbf{Keywords:} [关键词1]; [关键词2]; [关键词3]; [关键词4]

\noindent\textbf{JEL Codes:} [C23]; [D83]; [L86]
\end{abstract}

\newpage

% ============================================================
\section{Introduction}
% ============================================================

[引言正文。1500-2000词，5-7段。]

[P1: 研究背景与动机]

[P2-3: 文献空白]

[P4: 本文做了什么]

[P5: 核心发现]

[P6: 贡献]

The remainder of the paper is organized as follows. Section~\ref{sec:background} provides institutional background. Section~\ref{sec:literature} reviews related literature. Section~\ref{sec:data} describes the data. Section~\ref{sec:strategy} presents the empirical strategy. Section~\ref{sec:results} reports the results. Section~\ref{sec:conclusion} concludes.


% ============================================================
\section{Institutional Background}\label{sec:background}
% ============================================================

[制度/政策背景。1000-1500词。]


% ============================================================
\section{Literature Review}\label{sec:literature}
% ============================================================

[文献综述。1500-2000词，按主题组织。]


% ============================================================
\section{Data and Variables}\label{sec:data}
% ============================================================

\subsection{Data Sources}

[数据来源描述]

\subsection{Variable Definitions}

[变量定义]

\subsection{Descriptive Statistics}

Table~\ref{tab:summary} presents summary statistics for the main variables.

\begin{table}[H]
\centering
\caption{Summary Statistics}\label{tab:summary}
\begin{threeparttable}
\begin{tabular}{lccccc}
\toprule
Variable & N & Mean & SD & Min & Max \\
\midrule
[因变量] & [N] & [mean] & [sd] & [min] & [max] \\
[自变量] & [N] & [mean] & [sd] & [min] & [max] \\
[控制变量1] & [N] & [mean] & [sd] & [min] & [max] \\
[控制变量2] & [N] & [mean] & [sd] & [min] & [max] \\
\bottomrule
\end{tabular}
\begin{tablenotes}[flushleft]
\small
\item Notes: [说明数据来源、样本限制等]
\end{tablenotes}
\end{threeparttable}
\end{table}


% ============================================================
\section{Empirical Strategy}\label{sec:strategy}
% ============================================================

[识别策略说明]

The baseline specification is:
\begin{equation}\label{eq:did}
Y_{it} = \alpha_i + \delta_t + \beta \cdot \text{Treat}_i \times \text{Post}_t + X'_{it}\gamma + \varepsilon_{it}
\end{equation}
where $Y_{it}$ is [因变量] for unit $i$ in period $t$, $\alpha_i$ and $\delta_t$ are entity and time fixed effects, $\text{Treat}_i \times \text{Post}_t$ is the DID interaction, $X_{it}$ is a vector of controls, and $\varepsilon_{it}$ is the error term clustered at the [聚类层级] level.

The coefficient of interest is $\beta$, which captures [经济含义].


% ============================================================
\section{Results}\label{sec:results}
% ============================================================

\subsection{Main Results}

Table~\ref{tab:main} reports the main results.

\begin{table}[H]
\centering
\caption{Main Results}\label{tab:main}
\begin{threeparttable}
\begin{tabular}{lcccc}
\toprule
 & (1) & (2) & (3) & (4) \\
\midrule
Dep. Var.: & \multicolumn{4}{c}{[因变量]} \\
\midrule
Treat $\times$ Post & [coef]*** & [coef]*** & [coef]*** & [coef]*** \\
 & ([se]) & ([se]) & ([se]) & ([se]) \\
\midrule
Controls & No & Yes & Yes & Yes \\
Entity FE & Yes & Yes & Yes & Yes \\
Year FE & Yes & Yes & Yes & Yes \\
Province $\times$ Year FE & No & No & Yes & Yes \\
\midrule
Observations & [N] & [N] & [N] & [N] \\
$R^2$ (within) & [R2] & [R2] & [R2] & [R2] \\
\bottomrule
\end{tabular}
\begin{tablenotes}[flushleft]
\small
\item Notes: Robust standard errors clustered at the [level] level in parentheses. *** $p<0.01$, ** $p<0.05$, * $p<0.1$.
\end{tablenotes}
\end{threeparttable}
\end{table}

\subsection{Parallel Trends}

Figure~\ref{fig:event} presents the event study estimates.

\begin{figure}[H]
\centering
\includegraphics[width=0.85\textwidth]{figures/fig2_event_study.png}
\caption{Event Study: Dynamic Treatment Effects}\label{fig:event}
\begin{minipage}{0.85\textwidth}
\small
\textit{Notes:} [说明。基准期、置信区间、聚类方式等]
\end{minipage}
\end{figure}

\subsection{Robustness Checks}

[稳健性检验结果]

\subsection{Heterogeneity Analysis}

[异质性分析结果]


% ============================================================
\section{Conclusion}\label{sec:conclusion}
% ============================================================

[结论。800-1000词。核心发现→政策含义→局限→未来方向]


% ============================================================
% 参考文献
% ============================================================
\newpage
\bibliographystyle{aer}   % AER 格式; 或用 chicago, apa
\bibliography{references}  % references.bib 文件

% 如不用 bib 文件，可手动列出:
% \begin{thebibliography}{99}
% \bibitem{author2024} Author, A. (2024). Title. \textit{Journal}, 1(1), 1-20.
% \end{thebibliography}


% ============================================================
% 附录（如需要）
% ============================================================
\newpage
\appendix
\section{Additional Tables and Figures}\label{app:additional}

[附录表格和图]

\end{document}
